\section{The mechanism before the abstraction}
\label{sec:overview}

Suppose that a smooth projective variety $X$ determines finitely many local
blocks $a$, each with a positive multiplicity $m_X(a)$.  A comparison
theorem for a blowup does more than equate two unlabelled multisets: it matches
the blocks of the blowup with the ambient blocks and with a prescribed number
of copies of the center blocks.  Likewise, a projective-bundle comparison
matches its blocks with copies of the base blocks.  The elementary operation
laws become formal once each multiplicity is expanded into separate
occurrences and these comparison bijections are retained.

There are three successive constructions:
\[
 \text{local blocks with occurrences}
 \longrightarrow \text{thin comparison groupoid}
 \longrightarrow \text{free effective monoid}.
 \tag{1.1}
\]
The first arrow remembers exactly which occurrences the provider identifies.
The second forgets everything except their equivalence classes and
multiplicities.  A numerical marker is then simply a homomorphism out of the
free monoid.  Thus comparison is required before folding, but equality of
richer block objects after folding is not.

Katzarkov--Kontsevich--Pantev--Yu construct Hodge atoms from the Euler
spectral cover of a maximal non-archimedean A-model $F$-bundle
\cite[Sections~5.2--5.4]{KKPY}.  Their multiplicities are degrees of connected
components of the unramified reduced cover.  Their blowup and
projective-bundle decomposition theorems provide exactly the occurrence
correspondences used in (1.1).  The purpose of this note is to make the
resulting specialization and its universal property explicit.

The distinction between occurrences and isomorphism classes is essential.
It prevents a multiplicity from disappearing when two blocks are isomorphic,
and it separates two quotients that need not agree: the abstract atom quotient
generated by elementary geometric correspondences, and the quotient of
geometric atomic $F$-bundles by isomorphism.  The rationality criterion uses
the former.

Here is the theorem ladder.  Theorem~\ref{thm:ledger} constructs the universal
effective ledger from any occurrence provider.  Theorem~\ref{thm:hodge}
identifies the standard Hodge-atom chemical formula with its Hodge
specialization.  Theorem~\ref{thm:birational-quotient} kills precisely the
low-dimensional generators contributed by weak-factorization centers.
Corollary~\ref{cor:one-step} is the rank-two, one-step consequence.

The numerical direct-QDM atom used for cubic one-step irrationality in
\cite{Epilogue} is a sibling specialization of the same ledger mechanism,
not a Hodge atom under another name.  It uses a smaller block type and a
different fold.  No identification between the two carrier groupoids is
needed here.
